\subsection{\texorpdfstring{Equivalence classes of functions with values in \(\overline{R}\)}{Equivalence classes of functions with values in R-bar}}

In conformity with the definition in No.~3, we say that a function \(f\) , defined almost everywhere in \(X\) and with values in \(\overline{\mathbf{R}}\) , is finite almost everywhere if the set of \(x \in X\) for which \(f(x)\) is defined and finite has negligible complement. A function that is finite almost everywhere is equivalent to a function that is everywhere finite; one can therefore identify its class \(\widetilde{f}\) with a class of finite numerical functions defined on \(X\) (or almost everywhere in \(X\) ). In particular, the sum and product of two classes of functions finite almost everywhere are defined, and the set of these classes is an algebra over \(\mathbf{R}\) . If \((f_n)\) is a sequence of functions with values in \(\overline{\mathbf{R}}\) , defined and finite almost everywhere, then the partial sums \(\sum_{k=1}^{n} f_k(x)\) are defined almost everywhere; if, for almost every \(x \in X\) , they have a limit \(f(x)\) in \(\overline{R}\) , we again say that the series with general term \(f_{n}\) converges almost everywhere and that f is the sum of the series (note that f is not necessarily finite almost everywhere).

If \(f\) and \(g\) are two numerical functions defined and finite almost everywhere in \(X\) , then \(\widetilde{f} + \widetilde{g}\) (resp. \(\widetilde{f}\widetilde{g}\) ) is the class of every function equal to \(f(x) + g(x)\) (resp. \(f(x)g(x)\) ) at the points \(x \in X\) where this expression has meaning. Note that \(f\) and \(g\) can both be everywhere defined without \(f(x) + g(x)\) (resp. \(f(x)g(x)\) ) being defined for all \(x\) (GT, IV, §4, No.~3); by definition \(f + g\) (resp. \(fg\) ) is then the function equal to \(f(x) + g(x)\) (resp. \(f(x)g(x)\) ) at the points where this expression is defined; it is therefore only defined almost everywhere.

Let \(f\) and \(g\) be two numerical functions (finite or not) defined almost everywhere in \(X\) and such that \(f(x) \leqslant g(x)\) almost everywhere; if \(f_1\) is equivalent to \(f\) , and \(g_1\) is equivalent to \(g\) , it is clear that also \(f_1(x) \leqslant g_1(x)\) almost everywhere. The relation in question therefore depends only on the classes of \(f\) and \(g\) ; one writes \(\widetilde{f} \leqslant \widetilde{g}\) , and one verifies immediately that this relation is an order relation in the set of equivalence classes of functions with values in \(\overline{\mathbf{R}}\) . If \((\widetilde{f}_n)\) is a countable family (finite or infinite) of such classes and if, for every \(n\) , \(f_n\) and \(g_n\) are two functions defined almost everywhere and belonging to the class \(\widetilde{f}_n\) , it follows from Prop. 8 of No.~4 that the functions \(\sup_n f_n\) and \(\sup_n g_n\) , defined almost everywhere, are equivalent; their class therefore depends only on the classes \(\widetilde{f}_n\) , and one verifies at once that it is the supremum \(\sup_n \widetilde{f}_n\) of these classes in the set of classes of functions with values in \(\overline{\mathbf{R}}\) , ordered in the way just described (a set which is therefore, in particular, lattice-ordered). One shows similarly the existence of the infimum \(\inf_n \widetilde{f}_n\) , and one has \(\inf_n \widetilde{f}_n = -\sup_n (-\widetilde{f}_n)\) . It follows that \(\limsup_{n \to \infty} f_n\) and \(\limsup_{n \to \infty} g_n\) are also equivalent, and their class, which is denoted \(\limsup_{n \to \infty} \widetilde{f}_n\) , is equal to \(\inf_n (\sup_{p \geqslant 0} \widetilde{f}_{n+p})\) ; \(\liminf_{n \to \infty} \widetilde{f}_n\) is defined similarly.

A numerical function f (finite or not) is said to be negligible if it is equivalent to 0; this definition is equivalent to Def. 1 for functions that are positive and everywhere defined, by virtue of Th. 1. For f to be negligible, it is necessary and sufficient that \(|f|\) be negligible (or that both \(f^{+}\) and \(f^{-}\) be negligible).

